\chapter{文献综述}


\section{课堂提问的功能}

现阶段而言，国内外关于课堂提问功能的研究已基本成熟。

1967年，帕特选择小学教师作为研究对象，通过访谈进行分析。帕特的研究结果表明提问有五大功能：检查理解、辅助教学、诊断困难、刺激记忆和激发思维。这是首次对提问的功能进行总结。
1999年，施良方等人从原理、策略以及研究等方面阐述教学理论。书中提出提问的 4 种功能:诱发学生参与教学；提供线索；提供练习与反馈机会；有助于学生学习结果的迁移\cite{施良方1999教学理论}。 2002 年，加里在他的《有效教学方法》 中总结了提问的 6 种功能:吸引兴趣和吸引注意力；发现问题及检查；回忆具体知识或信息；课堂管理；鼓励更高层次的思维活动；组织或指导学习\cite{鲍里奇2002有效教学方法}。同年，2013 年， 朱杰珍归纳数学课堂提问的 5 种功能:温故 导学；激趣探疑；聚精促思；辨析明理；反馈调控\cite{朱杰珍2013刍议数学课堂提问的功能与技巧}。

可以看出，学者们都认为提问对教师和学生双方都是有利的，教师可以利用提问的教学手段帮助教学，学生也能由教师提问受到启发，从而进行更高层次的思维活动。



\section{课堂提问的类型}

课堂提问分类理论最早由巴恩斯于1969年提出。他将其分类为事实型问题、推理型问题、无需进行推 理的开放式问题、社交型问题。 而后，他在自己的研究基础上进行延伸，并发展出更加完整的理论体系，也就是将推理型问题分为封闭式与开放式。布鲁姆对提问分类提出了不同的观点，他认为可以按照问题的认知层次进行分类。依此标准，课堂提问划分为 6 种类型:记忆、理解、应用、分析、综合、评价。后续学者们的研究大多是以此为基础的。
对于小学数学课堂，我国学者王小清分析归纳 了四种课堂提问类型:猜测型提问；比较型提问；启发型提问；突破型提问\cite{王小清2020小学数学课堂提问类型探究}。2017 年，喻洋也对中学数学课堂提问类型进行分类:引入型问题；疏导型问题；探究型问 题；总结型问题；感想型问题\cite{喻洋2017中学数学课堂提问的类型与技巧}。

上述文献对课堂提问的分类进行研究，通常是根据提问的功能或者提问的方式两种视角进行研究，在加深对提问的认识上起到了关键作用，熟悉不同提问类型的特点有助于教师根据不同目标选取合适的提问类型。



\section{课堂提问的原则}

梁平基于课堂教学的效率、效益、体验这三大指标，针对目前初中数学课堂的提问现状，提出了课堂提问九大原则。该九大原则为：目的性原则——精心设计； 科学性原则——难易适度；趣味性原则——新颖别致；灵活性原则——因势利导；启发性原则——循循善诱；鼓励性原则——正确评价；广泛性原则——面向全体；针对性原则——因材施教；适量性原则——适量适度\cite{梁平2011初中数学课堂提问有效性及其策略的研究}。

黄小安结合结合阎承利学者提出的课堂提问的一般性原则，进一步提出具有高中数学学科特征的课堂提问原则：目的性原则；科学性原则；趣味性原则；启发性原则；灵活性和广泛性原则；鼓励性和针对性原则\cite{黄小安2006高中数学课堂提问有效性研究}。



\section{课堂提问的有效性}

随着对课堂提问研究的不断深入，学者们开始提出问题：如何检验课堂提问这一教学手段是否真正发挥其功能与作用？由此，关于课堂提问的有效性的研究逐渐受到关注。但是目前国内还未有统一的对有效课堂提问的认识。

2003 年，赵敏霞通过课堂观察及文献研究发现目前课堂提问现状不如人意，由 此她认为有效课堂提问应具有 3 个特点:是师生交流的重要形式；是实现教学目标整合的重要手段；是服务于学生的学习过程的\cite{赵敏霞2003对教师有效课堂教学提问的思考}。
2006年，王雪梅提出有效提问指教师清楚明了地、有目的、有组织地提出简短的、发人深省的问题，要能引起学生的回应和回答\cite{王雪梅2006课堂提问的有效性及其策略研究}。2010 年，卢正芝提出有效课堂提问的内涵。广义的有效课堂提问是指教师在精心预设问题的基础上，通过创设良好的问题情境，在教学中生成适切的问题，引导学生主动思考，进行质疑和对话，全面实现预期教学目标，并对提问及时反思与实践的过程\cite{卢正芝2010价值取向与标准建构}。
2011 年，温建红在《数学课堂有效提问的内涵及特征》中提出了有效 提问的几大特征:目的性；启发性；多样性；方法性；示范性；情感性\cite{温建红2011数学课堂有效提问的内涵及特征}。

综合以上种种研究，大家探讨课堂提问时大多针对普通教育而言，以数学教学为研究对象的研究比较少，而对复习课这一数学课型的课堂提问研究更是少之又少。因此，我希望对复习课这种数学课型进行探究，进一步分析在特定的数学课型当中课堂提问应遵循何种原则，希望对目前数学教学提供一定帮助。